\selectlanguage{ngerman}%
\frenchspacing%
\begin{ledgroupsized}[r]{120mm}%
\footnotesize%
\pstart%
\noindent\textbf{Überlieferung:}%
\pend%
\end{ledgroupsized}%
\begin{ledgroupsized}[r]{114mm}%
\footnotesize%
\pstart%
\parindent -6mm%
\makebox[6mm][l]{\textit{L}}%
Aufzeichnung:
LH XXXV 14,\;2 Bl.\;6.
Ein Blatt 4\textsuperscript{o},
Ränder beschnitten;
ein Wasserzeichen: %
römisches\indexO{Rom} Papier.
Zwei volle Seiten; %
Textfolge (editorisch festgelegt): Bl.\;\href{\Scan{LH_35_14_02_006v}}{6\,v\textsuperscript{o}}, \href{\Scan{LH_35_14_02_006r}}{6\,r\textsuperscript{o}}.
\pend%
\end{ledgroupsized}%
\vspace{5mm}%
\begin{ledgroup}%
\footnotesize%
\pstart%
\noindent%
\textbf{Datierungsgründe:}
Die Datierung der Aufzeichnung N.\;\ref{58220} beruht auf Eigenschaften des Textträgers sowie auf inhaltlichen Betrachtungen.
Siehe die ausführliche Begründung auf S.\;\pageref{58220_dat}.%
\pend%
\end{ledgroup}%
\selectlanguage{latin}%
\frenchspacing%
\vspace{7mm}%
\normalsize%
\pstart%
\noindent%
\edtext{%
\bracks{%
\href{\Scan{LH_35_14_02_006v}}%
{6\,v\textsuperscript{o}}%
}
Summum principium meum est%
}{%
\lemma{\bracks{\href{\Scan{LH_35_14_02_006v}}{6\,v\textsuperscript{o}}}}%
\Bfootnote{%
\textit{(1)}~Dico %
\textit{(2)}~Summum principium meum est%
~\textit{L}%
}}
in Mechanicis%
\indexS{principium summum in mechanicis}%
\indexS{mechanica}
effectum%
\indexS{effectus integrus}
\edtext{%
integrum%
}{%
\lemma{integrum}%
\Bfootnote{%
\textit{erg.\;L}%
}}
semper causae aequalem esse.%
\indexS{effectus aequalis causae}%
\indexS{causa aequalis effectui}
Unde consequitur
non tantum eandem semper servari summam virium,%
\indexS{summa virium servata}
sed et semper eandem summam directionum.%
\indexS{summa directionum servata}
\pend%
\pstart%
\edlabel{LH_35_14_02_006_006r_dmi-1}
Si idem punctum simul habeat diversas directiones,%
\indexS{punctum habens diversas directiones}%
\indexS{directiones puncti diversae}
ibit directione composita ex omnibus,%
\indexS{directio composita}
seu quod eodem
\edlabel{LH_35_14_02_006_006v3}%
redit,%
\edtext{%
}{%
\xxref{LH_35_14_02_006_006v3}{LH_35_14_02_006_006v4}%
\lemma{redit,}%
\Bfootnote{%
\textit{(1)}~perveniet ad centrum gravitatis dup %
\textit{(2)}~ibit %
\textit{(a)}~celeritate %
\textit{(b)}~directione centri gravitatis %
\textit{(aa)}~duplicationi %
\textit{(bb)}~suae duplicationis ad utrique nec %
\textit{(cc)}~quae foret %
\bracks{...} satisfaciendum necessariae, %
\textit{(aaa)}~per numerum %
\textit{(aaaa)}~multipl %
\textit{(bbbb)}~direction\angles{is} %
\textit{(cccc)}~directionum divisa. Ut si tres %
\textit{(bbb)}~et celeritate %
\bracks{...} numerum multiplicationis. \bracks{/} %
\textit{(aaaa)}~Si duo puncta in motu posita subito connecti ponantur, ibunt via et celeritate centri gravitatis communis, quae fuerat ante connexionem.\indexS{puncta mobilia connexa}\indexS{centrum gravitatis commune}\indexS{centrum gravitatis punctorum mobilium}\indexS{connexio punctorum mobilium} \bracks{/} %
\textit{(aaaaa)}~Idem est %
\textit{(bbbbb)}~Ac proinde si duo puncta connexa inter se solicitentur diversis directionibus, ponatur utraque succedere, quasi altera abesset, et notetur \bracks{/} %
\textit{(bbbb)}~Si duo %
\textit{(aaaaa)}~centri %
\textit{(bbbbb)}~puncta in motu %
\bracks{...} quae ante.%
~\textit{L}%
}}
ibit directione centri gravitatis%
\indexS{directio centri gravitatis}%
\indexS{centrum gravitatis}
quae foret in hypothesi multiplicationis suae
ad omnibus directionibus satisfaciendum necessariae,%
\indexS{hypothesis multiplicationis directionis}%
\indexS{multiplicatio directionis}
et celeritate ejusdem centri gravitatis,%
\indexS{celeritas centri gravitatis}%
\indexS{centrum gravitatis}
divisa per numerum multiplicationis.%
\edlabel{LH_35_14_02_006_006r_dmi-2}
\pend%
\begin{edfigure}[1.750em][0.00em]
\includegraphics[scale=0.80]{%
gesamttex/edit_VIII-4/images/LH_35_14_02_006_d1_006v.pdf}
\edcaption[0.0em]{\bracks{\textit{Fig.\:1}}}
\end{edfigure}
\newpage%
\pstart%
Si duo puncta in motu posita
subito connecti ponantur,%
\indexS{puncta mobilia connexa}%
\indexS{connexio punctorum mobilium}
centrum gravitatis eorum ibit%
\indexS{centrum gravitatis punctorum mobilium}
ea celeritate et directione%
\indexS{celeritas centri gravitatis}%
\indexS{directio centri gravitatis}
quae fuit ante concursum,%
\indexS{concursus punctorum mobilium}
et praeterea circulabuntur ambo
circa centrum gravitatis%
\indexS{circulatio circa centrum gravitatis}
ea celeritate,%
\indexS{celeritas circulationis}
ut eadem sit summa virium%
\indexS{summa virium servata}
quae ante.%
\edlabel{LH_35_14_02_006_006v4}%
\pend%
\pstart%
Sint
\edtext{%
puncta \textit{A} et \textit{B},%
}{%
\lemma{puncta \textit{A} et \textit{B}}%
\Cfootnote{%
Vgl. das Diagramm \bracks{\textit{Fig.\:1}}.%
}}
unius velocitas et directio \textit{A(A)},
alterius \textit{B(B)}.%
\indexS{velocitas punctorum mobilium}%
\indexS{directio punctorum mobilium}
Connectuntur in \textit{(A)(B)}.%
\indexS{connexio punctorum mobilium}
Si ejusdem ponantur
\edtext{%
densitatis%
\indexS{densitas punctorum mobilium}
erit%
}{%
\lemma{densitatis}%
\Bfootnote{%
\textit{(1)}~seu sibi %
\textit{(2)}~erit%
~\textit{L}%
}}
centrum gravitatis%
\indexS{centrum gravitatis punctorum mobilium}
in
\edtext{%
medio ipsius \textit{AB}
et ipsius \textit{(A)(B)}.%
}{%
\lemma{medio}%
\Bfootnote{%
\textit{(1)}~, et %
\textit{(2)}~ipsarum \textit{A} %
\textit{(3)}~ipsius \textit{AB} et ipsius \textit{(A)(B)}.%
~\textit{L}%
}}
Ergo celeritas et directio centri
est \textit{C(C)},%
\indexS{celeritas centri gravitatis}%
\indexS{directio centri gravitatis}
cui aequalis debet esse \textit{(C)((C))},
cui si accedat circulatio,%
\indexS{circulatio circa centrum gravitatis}
debet esse talis %
ut sit summa quadratorum
$ \overline{\textit{(A)((A))}}^2 + \overline{\textit{(B)((B))}}^2 $
aequalis summae quadratorum%
\indexS{summa quadratorum velocitatum}%
\indexS{quadratum velocitatis}
$ \textit{A(A)}^2 $ et 
\edlabel{LH_35_14_02_006_006v5}%
$ \textit{B(B)}^2 $.%
\edtext{}{%
\xxref{LH_35_14_02_006_006v5}{LH_35_14_02_006_006v6}%
\lemma{$B(B)^2$.}%
\Bfootnote{%
\bracks{/}
\textit{(1)}~Id %
\textit{(2)}~Eodem %
\bars{igitur \textit{erg.}} modo%
~\textit{L}%
}}%
\pend%
\pstart%
Eodem igitur modo%
\edlabel{LH_35_14_02_006_006v6}
si duo puncta jam connexa%
\indexS{puncta mobilia connexa}
\textit{(A)} et \textit{(B)}
solicitentur
directionibus%
\indexS{directio punctorum mobilium}
\textit{A(A)} et \textit{B(B)},
directiones inde ortae erunt
\textit{(A)((A))} et \textit{(B)((B))}.%
\indexS{directio punctorum mobilium}
Videndum
an tunc contingat
nullam fieri circulationem,%
\indexS{circulatio circa centrum gravitatis}
quando vires%
\indexS{vis punctorum mobilium}
seu celeritates%
\indexS{celeritas punctorum mobilium}
(punctis positis aeque gravibus)%
\indexS{puncta mobilia aeque gravia}
sunt reciproce proportionales
distantiis a centro gravitatis \textit{C},%
\indexS{distantia a centro gravitatis}
ut tunc scilicet recta \textit{((A))((B))}
\edtext{%
man\bracks{et}%
}{%
\lemma{manent}%
\Bfootnote{%
\textit{L\;ändert Hrsg.}%
}}
parallela rectae \textit{(A)(B)}.
Videndum
an non res eodem redeat,
ac si ponas \textit{(C)}
\edtext{%
hypomochlium%
\indexS{hypomochlium}%
}{%
\lemma{hypomochlium}%
\Bfootnote{%
\textit{erg.\;L}%
}}
moveri motu uniformi,%
\indexS{motus uniformis hypomochlii}
et puncta
\edtext{%
\textit{A} et \textit{B} niti circulari
ut gravia aequalia posita%
\indexS{gravia aequalia}
in \textit{P} et \textit{E},
posito%
}{%
\lemma{\textit{A} et \textit{B}}%
\Bfootnote{%
\textit{(1)}~posita in %
\textit{(2)}~niti circulari %
\bracks{...} \textit{P} et \textit{E}, %
\textit{(a)}~\angles{--} %
\textit{(b)}~posito%
~\textit{L}%
}}
\textit{(C)P} et \textit{(C)E}
esse perpendiculares in directiones,
quanquam ea res
\edtext{%
sine nostris%
}{%
\lemma{sine nostris}%
\Cfootnote{%
Anspielung unklar.%
}}
ad determinandum
modum
non sufficiat.
\bracks{%
\href{\Scan{LH_35_14_02_006r}}%
{6\,r\textsuperscript{o}}%
}
\pend%
\begin{edfigure}[2.50em][0.00em]
\includegraphics[scale=1.0]%
{gesamttex/edit_VIII-4/images/LH_35_14_02_006_d2_006r.pdf}%
\edcaption{%
\hspace{1.0mm}%
\bracks{\textit{Fig.\:2}}%
}%
\end{edfigure}
\newpage%
\pstart%
\edlabel{LH_35_14_02_006_006r_mip-1}%
Quippe demonstravimus
\edtext{%
alias,%
}{%
\lemma{alias}%
\Cfootnote{%
Siehe die Erläuterung zu \refpassage{LH_35_14_02_006_006r_njl-1}{LH_35_14_02_006_006r_njl-2}.
}}
si tractiones in%
\indexS{tractio in extremis}%
\indexS{tractio circa hypomochlium}
\edtext{%
extremis\;\textit{A} et\;\textit{B}%
}{%
\lemma{extremis \textit{A} et \textit{B}}%
\Cfootnote{%
Vgl. das Diagramm \bracks{\textit{Fig.\:2}}.%
}}
circa
\edtext{%
hypomochlium\;\textit{C}%
\indexS{hypomochlium}
fiant,
vires circulandi%
\indexS{vis circulandi}
fore%
}{%
\lemma{hypomochlium\;\textit{C}}%
\Bfootnote{%
\textit{(1)}~sint %
\textit{(2)}~fiant, %
\textit{(a)}~vim fore %
\textit{(b)}~vires circulandi fore%
~\textit{L}%
}}
in composita ratione%
\indexS{ratio composita}
\edtext{%
tractionum
per \textit{EB} et \textit{DA},
et perpendicularium
\textit{CF}, \textit{CG}
ex Hypomochlio \textit{C}%
\indexS{hypomochlium}
in directiones \textit{AD}, \textit{BE}
ductarum.%
}{%
\lemma{tractionum}%
\Bfootnote{%
\textit{(1)}~, et %
\textit{(2)}~\textit{EB} \textit{DA}, %
\textit{(3)}~per \textit{EB} et \textit{DA}, et perpendi\-cularium %
\bars{\textit{CF}, \textit{CG} \textit{erg.}} ex Hypomochlio \textit{C} in %
\textit{(a)}~lineas %
\textit{(b)}~directiones \textit{AD}, %
\textit{(aa)}~\textit{AE}, du %
\textit{(bb)}~\textit{BE} ductarum.%
~\textit{L}%
}}
\edlabel{LH_35_14_02_006_006r_njl-1}%
\edtext{%
Videantur mea ad posthuma Fermatii.%
\indexN{\textls{Fermat}, Pierre de 1601\textendash1665}%
\indexS{posthumum}%
}{%
\lemma{Videantur \bracks{...} Fermatii}%
\Cfootnote{%
Zu \textsc{P.\;de Fermat},
\textit{Varia opera mathematica} %
(Toulouse 1679)\cite{00533}
hat Leibniz ausführlich kommentierte Auszüge verfasst;
vgl. N.\;\ref{58229} im vorliegenden Band. %
Siehe auch die ausführlichen Datierunggründe zu N.\;\ref{58220}, S.\;\pageref{58220_dat}.%
}}%
\edlabel{LH_35_14_02_006_006r_njl-2}
Tractiones autem aestimandae%
\indexS{tractio aestimanda}
in ratione composita%
\indexS{ratio composita}
celeritatum et ponderum,%
\indexS{ratio celeritatum et poderum}
quia hic vis mortuae ratio%
\indexS{vis mortua}%
\indexS{ratio vis mortuae}
\edlabel{LH_35_14_02_006_006r1}%
habenda.%
\edlabel{LH_35_14_02_006_006r_mip-2}%
\edtext{}{%
\xxref{LH_35_14_02_006_006r1}{LH_35_14_02_006_006r2}%
\lemma{habenda.}%
\Bfootnote{%
\bracks{/}
\textit{(1)}~Quod si ea methodus succedat, ut liceat aestimare motum ex motu centri et circulatione circa centrum quasi hypomochlium,\indexS{aestimatio motus} %
\textit{(a)}~licebit jam de integro aliquo corpore sic ratiocinari. %
\textit{(aa)}~Sint %
\textit{(bb)}~Sit corpus %
\textit{(b)}~seu ut %
\textit{(aa)}~sumatur %
\textit{(bb)}~sumantur corpora %
\textit{(c)}~et plurium %
\textit{(aa)}~corporum %
\textit{(bb)}~punctorum \textit{A, B, D} centrum gravitatis commune %
\bars{\textit{C} \textit{erg.}} sumere, et %
\bars{ita \textit{gestr.}} fingere %
\bracks{/}
\textit{(2)}~Si plura %
\bracks{...} an res %
\textit{(a)}~\angles{ex} %
\textit{(b)}~nostra methodo %
\bracks{...} tria puncta %
\bars{\textit{A, B, C} connexa inter se et \textit{erg.}} conantia %
\textit{(aa)}~velocita %
\textit{(bb)}~aequalibus %
\textit{(cc)}~tractionibus %
\textit{(aaa)}~et %
\textit{(bbb)}~\textit{AL}, \textit{BM}, \textit{CN}, et quaeratur primum %
\textit{(aaaa)}~via %
\textit{(bbbb)}~directio et %
\bracks{...} ut summae %
\textit{(aaaaa)}~velocitatum %
\textit{(bbbbb)}~quadratorum a %
\bracks{...} directione progrederetur;%
~\textit{L}%
}}%
\pend%
\pstart%
Si plura sint puncta,
videndum an res
\edtext{%
nostra methodo prima%
\indexS{methodus nostra}%
}{%
\lemma{nostra methodo prima}%
\Cfootnote{%
Vgl. oben, \refpassage{LH_35_14_02_006_006r_dmi-1}{LH_35_14_02_006_006r_dmi-2}.%
}}
ex sola virium et directionum consideratione%
\indexS{consideratio virium et directionum}
sit determinata.
Sint
\edtext{%
tria puncta \textit{A,\;B,\;C}%
}{%
\lemma{tria puncta \textit{A,\;B,\;C}}%
\Cfootnote{%
Vgl. das Diagramm \bracks{\textit{Fig.\:3}}.%
}}
connexa inter se%
\indexS{puncta mobilia connexa}
et conantia tractionibus%
\indexS{punctum conans tractione}
\textit{AL}, \textit{BM}, \textit{CN},
et quaeratur primum
directio et velocitas centri gravitatis,%
\indexS{directio centri gravitatis}%
\indexS{velocitas centri gravitatis}%
\indexS{centrum gravitatis punctorum mobilium}
deinde talis circulatio%
\indexS{circulatio circa centrum gravitatis}
ut summae quadratorum a velocitatibus%
\indexS{summa quadratorum velocitatum}%
\indexS{quadratum velocitatis}
in ipsorum punctorum massas ductorum%
\indexS{massa punctorum mobilium}
sint eaedem
quae forent %
si quodlibet punctum sua directione progrederetur;%
\indexS{directio punctorum mobilium}%
\edlabel{LH_35_14_02_006_006r2}
et habebitur
\edlabel{LH_35_14_02_006_006r3}%
solutio.%
\indexS{solutio}%
\edtext{}{%
\xxref{LH_35_14_02_006_006r3}{LH_35_14_02_006_006r4}%
\lemma{solutio.}%
\Bfootnote{%
\bars{Hoc tamen %
\bracks{...} solicitationibus solius %
\textit{(1)}~velocitatibus %
\textit{(2)}~velocitatis non %
\bracks{...} ratio habeatur. \textit{erg.}} %
\bracks{/} %
\textit{(1)}~Videndum an non %
\textit{(2)}~Et videndum an non%
~\textit{L}%
}}
Hoc tamen notandum,
ut in primis seu infinite parvis solicitationibus%
\indexS{solicitatio prima seu infinite parva}
solius velocitatis%
\indexS{solicitatio velocitatis}%
\indexS{velocitas solicitata}%
\bracks{,}
non quadratorum ejus%
\bracks{,}
ratio habeatur.
\pend%
\pstart%
Et videndum
an non%
\edlabel{LH_35_14_02_006_006r4}
idem obtineatur
\edtext{%
methodo illa priore,%
\indexS{methodus prior}%
}{%
\lemma{methodo illa priore}%
\Cfootnote{%
Vgl. oben \refpassage{LH_35_14_02_006_006r_mip-1}{LH_35_14_02_006_006r_mip-2}.%
}}
ut dato motu debito ipsi centro%
\indexS{motus centri gravitatis}%
\indexS{centrum gravitatis}
\edtext{%
gravitatis%
\bracks{,}
conatus circulandi%
\indexS{conatus circulandi}%
}{%
\lemma{gravitatis}%
\Bfootnote{%
\textit{(1)}~vires ci %
\textit{(2)}~conatus circulandi%
~\textit{L}%
}}
ponantur esse ut normales in lineas directionum.%
\indexS{linea directionis}
Etsi id quidem non sufficiat
ad absolute habendam circulationis%
\indexS{circulatio circa centrum gravitatis}
\edlabel{LH_35_14_02_006_006r5}%
velocitatem.%
\indexS{velocitas circulationis}%
\edtext{}{%
\xxref{LH_35_14_02_006_006r5}{LH_35_14_02_006_006r6}%
\lemma{velocitatem.}%
\Bfootnote{%
\bracks{/} %
\textit{(1)}~Si ponamus pu %
\textit{(2)}~Si puncta ponantur%
~\textit{L}%
}}%
\pend%
\newpage%
\pstart%
Si puncta ponantur%
\edlabel{LH_35_14_02_006_006r6}
esse infinita%
\indexS{puncta mobilia infinita}
nihilominus eorum centrum gravitatis%
\indexS{centrum gravitatis punctorum infinitorum}
serviet
et circulatio determinata habebitur
ex conservatione virium,%
\indexS{conservatio virium}
vel ex proportione normalium ad lineas%
\indexS{linea directionis}
\edtext{%
directionum.
Tale quid in
\edtext{%
linea \textit{AB}:%
}{%
\lemma{linea \textit{AB}}%
\Cfootnote{%
Vgl. das Diagramm \bracks{\textit{Fig.\:4}}.%
}}
si integra esset figura plana%
\indexS{figura plana integra}%
\bracks{,}
eadem esset methodus,%
\indexS{methodus}%
}{%
\lemma{directionum.}%
\Bfootnote{%
\textit{(1)}~Si integra sit figura %
\textit{(2)}~Tale quid %
\bracks{...} esset figura %
\textit{(a)}~ead %
\textit{(b)}~plana eadem esset methodus,%
~\textit{L}%
}}
datur tamen compendium%
\indexS{compendium}
si omnes lineae directionum sint eodem plano,
nam punctum \textit{D} intelligitur grave%
\indexS{punctum mobile grave}
in ratione rectae \textit{FD},
quia cadit in lineam directionis%
\indexS{linea directionis}%
\lbrack,
\edtext{%
et\rbrack%
}{%
\lemma{in}%
\Bfootnote{%
\textit{L\;ändert Hrsg.}%
}}%
\;\textit{E} in ratione rectae \textit{EH}.
\pend%
\pstart%
Data
\edtext{%
\edtext{%
figura%
}{%
\lemma{\textit{Über dem Wort} figura \textit{vermerkt:}}%
\Afootnote{%
positione%
}}
aliqua \textit{ABC}%
}{%
\lemma{figura aliqua \textit{ABC}}%
\Cfootnote{%
Vgl. das Diagramm \bracks{\textit{Fig.\:5}}.%
}}
ad centrum terrae \textit{G} tendente,%
\indexS{tendentia ad centrum terrae}%
\indexS{centrum terrae}
invenire lineam descensus%
\indexS{linea descensus}%
\indexS{descensus punctorum mobilium}
seu directionis totalis,%
\indexS{linea directionis totalis}%
\indexS{directio totalis punctorum mobilium}
sive punctum sustentationis.%
\indexS{punctum sustentationis}%
\indexS{sustentatio}
\pend%
\begin{edfigure}[3.50em][2.00em]
\hfill%
\includegraphics[scale=1.0]{%
gesamttex/edit_VIII-4/images/LH_35_14_02_006_d3_006r.pdf}
\hfill%
\hspace*{1mm}%
\includegraphics[scale=1.0]{%
gesamttex/edit_VIII-4/images/LH_35_14_02_006_d4_006r.pdf}
\hspace*{4mm}%
\hfill%
\includegraphics[scale=1.0]{%
gesamttex/edit_VIII-4/images/LH_35_14_02_006_d5_006r.pdf}
\hfill%
\edcaption[0.75em]{%
\hspace*{18.5mm}%
\bracks{\textit{Fig.\:3}}%
\hspace*{30mm}%
\bracks{\textit{Fig.\:4}}%
\hspace*{34,5mm}%
\bracks{\textit{Fig.\:5}}%
\hfill%
\edtext{}{%
\lemma{\bracks{\textit{Fig.\:3}} bis \bracks{\textit{Fig.\:5}}}%
\Cfootnote[nonum]{%
\textit{L} überliefert (Bl.\;\href{\Scan{LH_35_14_02_006r}}{6\,r\textsuperscript{o}}) gestrichene, hier nicht wiedergegebene Entwürfe zu den Diagrammen.%
}}
}%
\end{edfigure}
\vspace{3.0em}%
\pstart%
\noindent%
\centering%
*     *     *
\pend%
\newpage%
\selectlanguage{german}%
\frenchspacing%
\footnotesize%
\pstart%
\noindent%
\textls{Ausführliche Datierungsgründe:}%
\pagelabel{58220_dat}
Der Träger der titellosen und undatierten Aufzeichnung N.\;\ref{58220} weist ein Wasserzeichen (umkreiste sechsstrahlige Sonne mit Kreuz und Buchstaben F) auf, welches die römische\indexO{Rom} Herkunft des Papiers bezeugt.
Auf gleichem Papier hat Leibniz zahlereiche weitere Texte verfasst, die allesamt auf die Zeit seines Aufenthalts in Rom\indexO{Rom} -- Mitte April bis November 1689 -- datierbar und in diesem Band ediert sind: u.\,a. N.\;\ref{58656-A}, N.\;\ref{58258}, N.\;\ref{58218}, N.\;\ref{58240}, N.\;\ref{60197} und N.\;\ref{58235} (siehe für die Details die jeweiligen Datierungsgründe, vornehmlich die zu N.\;\ref{58656-A}, S.\;\pageref{58656-A_dat}). %
Besonders relevant ist, dass auf Papier dieser Art auch Leibnizens ausführliche Auzüge mit Bemerkungen aus \textsc{P.\;de Fermat}, \textit{Varia opera mathematica} %
(Toulouse 1679)\cite{00533} überliefert sind, auf die in der Aufzeichnung N.\;\ref{58220} (\refpassage{LH_35_14_02_006_006r_njl-1}{LH_35_14_02_006_006r_njl-2}) ausdrücklich hingewiesen wird;
Leibniz beieht sich hierbei auf Ergebnisse seiner Auseinandersetzung mit Fragen der Geostatik, die er vorwiegend im Briefwechsel zwischen Pierre de Fermat\indexN{\textls{Fermat}, Pierre de 1601\textendash1665} und Gilles Personne der Roberval\indexN{\textls{Roberval} (Robervallius), Gilles Personne de 1602\textendash1675} lesen konnte (vgl. N.\;\ref{58229} im vorliegenden Band).
Foglich muss N.\;\ref{58220} nach den kommentierten Auszügen N.\;\ref{58229} abgefasst worden sein.
Eine Entstehung der Aufzeichnung N.\;\ref{58220} vor Mitte April 1689 ist somit ausgeschlossen;
eine Entstehung nach November 1689 unwahrscheinlich.
\pend%
\selectlanguage{latin}%
\frenchspacing%
\normalsize%
